%% Section 1 -- Introduction & Motivation
%% Ported from docs/DRAFT_PAPER.md Section 1. British spelling is deliberate.

\section{Introduction \& Motivation}
\label{sec:introduction}

A phone stops working and goes in for repair. Its owner keeps the messaging
account running on a spare handset, and the conversations continue there for as
long as the repair takes. The original phone comes back. The owner restores the
most recent backup onto it, and the messages from the repair window vanish from
the screen.

Nothing was deleted. Those messages still sit in the spare device's local
database, and in any export the user thought to take before handing the account
back. The platform simply offers no operation that puts the two histories
together. WhatsApp's own restore guidance requires the same account that created
the backup~\citep{whatsapp_backup_restore}; it documents no operation for
combining two backups, a gap that has motivated the third-party merge tools
surveyed in Section~\ref{sec:related:tooling}.

We call this state \textbf{Temporal Communication Fragmentation} (TCF). Its
signature is a fully intact archive that the platform will not surface. The data
survives; the affordance to reach it does not.

Three properties make the situation harder than it first appears. The two
archives overlap in part and diverge in part, so neither is a superset of the
other. They carry no shared version history, no operation log, and no common
ancestor, because consumer chat exports were never designed to be reconciled.
And the restore operation that finalises the split gives no warning that a
competing archive exists, attempts no comparison, and reports nothing about what
it replaced.

\paragraph{Motivation, stated as motivation.}
This work began with one of the authors losing a stretch of message history to
exactly this sequence: primary handset unavailable, account carried on a second
handset, history replaced on restore. That incident prompted the research
questions below. It is an anecdote. It is not evidence, no controlled experiment
was run against it, and no claim in this paper rests on it.

The paper asks three questions:

\begin{description}
  \item[RQ1.] What conditions must hold for a single account's message history to
    end up split across two archives that are each internally complete and
    mutually irreconcilable through the platform?
  \item[RQ2.] Given two or more such archives, what does a reconstruction
    procedure have to guarantee for its output to be trustworthy as a record,
    rather than merely plausible?
  \item[RQ3.] Where do those guarantees sit relative to what optimistic
    replication, digital forensics, and the existing body of chat-merging tools
    already provide?
\end{description}

Answering RQ1 gives a formal characterisation of TCF with three necessary
conditions (Section~\ref{sec:tcf}). Answering RQ2 gives CARTA (Chat Archive
Reconstruction \& Temporal Access), a platform-agnostic specification of a
three-stage reconstruction contract, \carta, in which the data passed between
stages is typed, records that cannot be safely interpreted are retained in a
quarantine set $\mathcal{Q}$ rather than dropped, and every reconstruction emits
a structured report $\Omega$ covering provenance, quarantine carryover,
conflicts, and counts (Sections~\ref{sec:carta} and~\ref{sec:algorithms}).
Answering RQ3 is Section~\ref{sec:related}.

\paragraph{What this paper does not do.}
No implementation is evaluated here. CHAMP-Pipeline, the WhatsApp-specific
reference implementation of CARTA, is described in
Section~\ref{sec:implementation} as planned work and does not yet exist;
\texttt{core/src/\{ingestion,normalization,synthesis\}/} are empty. There is
accordingly no benchmark, no performance measurement, and no
reconstruction-accuracy figure anywhere in this paper. Where a capability is
stated, it is either a property the specification requires of a conforming
implementation or a citation to someone else's system. Readers should treat the
contribution as a problem definition and a specification, with implementation
and evaluation left as declared future work.
